% Chapitre 1 : Entités à annoter et leurs relations

\section{Entités à annoter et leurs relations}

\subsection{Vue d'ensemble des catégories d'annotation}

Le schéma d'annotation pour l'identification des bactériémies comprend six catégories principales d'entités et quatre types de relations. Chaque entité doit être annotée uniquement lorsqu'elle est directement liée à un épisode de bactériémie confirmé.

\begin{table}[h]
\centering
\caption{Synthèse des entités à annoter}
\begin{tabular}{|l|p{8cm}|l|}
\hline
\textbf{Entité} & \textbf{Description} & \textbf{Couleur BRAT} \\
\hline
BACTÉRIÉMIE & Mention de l'infection bactérienne dans le sang & Rouge \\
BACTÉRIE & Agent pathogène responsable & Bleu \\
SITE\_PRIMAIRE & Porte d'entrée de l'infection & Vert \\
SITE\_SECONDAIRE & Localisation des complications & Orange \\
INFECTION & Description du foyer infectieux & Violet \\
RÉSISTANCE & Profil de résistance antibiotique & Jaune \\
\hline
\end{tabular}
\end{table}

\subsection{Entité BACTÉRIÉMIE}

\subsubsection{Définition}
L'entité BACTÉRIÉMIE identifie toute mention explicite ou implicite d'une infection bactérienne dans le sang confirmée.

\subsubsection{Termes à annoter}
\begin{itemize}
    \item \textbf{Termes directs} : bactériémie, septicémie, sepsis, choc septique
    \item \textbf{Termes de laboratoire} : hémoculture positive, hémocultures positives, bactériémiant
    \item \textbf{Formulations cliniques} : "germe dans le sang", "infection du sang", "passage sanguin"
\end{itemize}

\subsubsection{Exemples d'annotation}
\begin{verbatim}
✓ "Le patient présente une [bactériémie] à Staphylocoque doré"
✓ "Les [hémocultures] sont revenues positives à E. coli"
✓ "Diagnostic de [septicémie] sur cathéter central"
✗ "Suspicion de bactériémie" (non confirmée)
✗ "Hémocultures négatives" (absence de bactériémie)
\end{verbatim}

\subsubsection{Règles spécifiques}
\begin{enumerate}
    \item Annoter uniquement les bactériémies confirmées microbiologiquement
    \item Inclure l'article et les modificateurs dans l'annotation ("une bactériémie", "la septicémie")
    \item Ne pas annoter les suspicions, hypothèses ou bactériémies écartées
\end{enumerate}

\subsection{Entité BACTÉRIE}

\subsubsection{Définition}
L'entité BACTÉRIE identifie l'agent pathogène responsable de la bactériémie, incluant toutes les variantes nomenclaturales.

\subsubsection{Catégories de bactéries à annoter}

\paragraph{Cocci Gram positif}
\begin{itemize}
    \item \textbf{Staphylocoques} : Staph, Staph doré, Staph blanc, S. aureus, Staphylococcus aureus, SARM, SASM, Staphylocoque coagulase négative, S. epidermidis, S. hominis, S. haemolyticus
    \item \textbf{Streptocoques} : Strepto, Streptocoque, S. pneumoniae, Pneumocoque, Streptocoque du groupe A/B/C/D, S. pyogenes, S. agalactiae, S. dysgalactiae, Enterocoque, E. faecalis, E. faecium, ERV
    \item \textbf{Autres} : Micrococcus, Pediococcus, Aerococcus, Granulicatella
\end{itemize}

\paragraph{Bacilles Gram négatif}
\begin{itemize}
    \item \textbf{Entérobactéries} : E. coli, Escherichia coli, Klebsiella, K. pneumoniae, Enterobacter, Citrobacter, Proteus, P. mirabilis, Serratia, S. marcescens, Salmonella, Shigella, Morganella, M. morganii, Providencia, Hafnia, Pantoea, Raoultella
    \item \textbf{Non fermentants} : Pseudomonas, P. aeruginosa, Acinetobacter, A. baumannii, Stenotrophomonas, S. maltophilia, Burkholderia, Achromobacter, Elizabethkingia, Chryseobacterium
    \item \textbf{Autres BGN} : Haemophilus, H. influenzae, Moraxella, Neisseria, Pasteurella, Aeromonas, Vibrio, Campylobacter, Helicobacter
\end{itemize}

\paragraph{Anaérobies}
\begin{itemize}
    \item Clostridium, C. perfringens, C. difficile
    \item Bacteroides, B. fragilis
    \item Fusobacterium, Prevotella, Propionibacterium
    \item Peptostreptococcus, Anaerococcus, Parvimonas
\end{itemize}

\paragraph{Autres}
\begin{itemize}
    \item \textbf{Groupe HACEK} : Haemophilus, Aggregatibacter, Cardiobacterium, Eikenella, Kingella
    \item \textbf{Intracellulaires} : Legionella, Brucella, Francisella, Bordetella
    \item \textbf{Filamenteux} : Nocardia, Actinomyces
    \item \textbf{Mycobactéries} : M. tuberculosis, mycobactéries atypiques
    \item \textbf{Levures} : Candida, C. albicans, C. glabrata (bien que techniquement fongémie)
\end{itemize}

\subsubsection{Exemples d'annotation}
\begin{verbatim}
✓ "Bactériémie à [Staphylococcus aureus] résistant à la méthicilline"
✓ "Présence de [E. coli] BLSE dans les hémocultures"
✓ "Septicémie à [BGN] non identifié"
✓ "[Streptocoque du groupe B] isolé sur 3 flacons"
\end{verbatim}

\subsection{Entité SITE\_PRIMAIRE}

\subsubsection{Définition}
Le SITE\_PRIMAIRE identifie la porte d'entrée initiale de l'infection bactérienne dans l'organisme.

\subsubsection{Catégories de sites primaires}

\begin{table}[h]
\centering
\caption{Sites primaires et termes associés}
\begin{tabular}{|l|p{10cm}|}
\hline
\textbf{Site} & \textbf{Termes et expressions associés} \\
\hline
Urinaire & Urines, urinaire, pyélonéphrite, PNA, cystite, prostatite, infection prostatique, sonde urinaire, sonde vésicale, foyer rénal \\
\hline
Digestif & Digestif, biliaire, cholangite, cholécystite, péritonite, abcès hépatique, diverticulite, appendicite, colite, entérite, translocation digestive \\
\hline
Pulmonaire & Poumon, pulmonaire, pneumonie, pneumopathie, pleuropneumopathie, abcès pulmonaire, empyème \\
\hline
Cathéter & Cathéter, KT, voie veineuse centrale, VVC, PAC, chambre implantable, PICC-line, cathéter de dialyse, infection sur cathéter \\
\hline
Cutané & Peau, cutané, érysipèle, cellulite, dermohypodermite, escarre, plaie, abcès cutané, fasciite, gangrène \\
\hline
Ostéo-articulaire & Ostéo-articulaire, ostéite, ostéomyélite, arthrite, spondylodiscite, infection sur prothèse, infection osseuse \\
\hline
Cardio-vasculaire & Endocardite, péricardite, myocardite, infection sur pacemaker, infection sur valve, greffe vasculaire infectée \\
\hline
ORL/Stomato & ORL, sinusite, otite, mastoïdite, pharyngite, angine, abcès dentaire, cellulite cervico-faciale \\
\hline
Génital & Génital, endométrite, salpingite, abcès tubo-ovarien, orchite, épididymite \\
\hline
SNC & Méningite, méningoencéphalite, abcès cérébral, empyème sous-dural \\
\hline
\end{tabular}
\end{table}

\subsubsection{Expressions types}
\begin{itemize}
    \item "point de départ [site]"
    \item "porte d'entrée [site]"
    \item "sur foyer [site]"
    \item "origine [site]"
    \item "secondaire à une infection [site]"
    \item "dans un contexte de [infection site]"
\end{itemize}

\subsection{Entité SITE\_SECONDAIRE}

\subsubsection{Définition}
Le SITE\_SECONDAIRE identifie les localisations de dissémination ou de complications de la bactériémie.

\subsubsection{Termes spécifiques}
\begin{itemize}
    \item \textbf{Termes de dissémination} : embole septique, greffe septique, métastase septique, localisation secondaire
    \item \textbf{Expressions types} : "compliquée de", "avec atteinte", "extension vers", "dissémination au niveau"
\end{itemize}

\subsubsection{Exemples d'annotation}
\begin{verbatim}
✓ "Bactériémie compliquée d'[emboles septiques pulmonaires]"
✓ "Avec [greffe septique sur la valve mitrale]"
✓ "Extension avec [spondylodiscite L4-L5]"
\end{verbatim}

\subsection{Entité INFECTION}

\subsubsection{Définition}
L'entité INFECTION décrit le type spécifique d'infection associée à la bactériémie, qu'elle soit primaire ou secondaire.

\subsubsection{Exemples par système}
\begin{itemize}
    \item \textbf{Urinaire} : pyélonéphrite aiguë, PNA, abcès rénal, pyonéphrose
    \item \textbf{Digestif} : péritonite, cholangite, abcès hépatique, appendicite compliquée
    \item \textbf{Pulmonaire} : pneumonie lobaire, abcès pulmonaire, empyème pleural
    \item \textbf{Cutané} : érysipèle, fasciite nécrosante, gangrène gazeuse
    \item \textbf{Ostéo-articulaire} : ostéomyélite aiguë, arthrite septique, spondylodiscite
\end{itemize}

\subsection{Entité RÉSISTANCE}

\subsubsection{Définition}
L'entité RÉSISTANCE identifie le profil de résistance ou de sensibilité aux antibiotiques de la bactérie.

\subsubsection{Termes à annoter}
\begin{itemize}
    \item \textbf{Résistances spécifiques} : SARM, BLSE, EPC, VRE, résistant à la méthicilline, résistant aux fluoroquinolones
    \item \textbf{Sensibilités} : sensible, sauvage, SASM, pan-sensible
    \item \textbf{Marqueurs} : PLP2a positif, carbapénémase, bêta-lactamase
\end{itemize}

\subsection{Relations entre entités}

\subsubsection{Relation CAUSÉE\_PAR}
\begin{itemize}
    \item \textbf{Type} : Bactériémie → Bactérie
    \item \textbf{Cardinalité} : 1 bactériémie peut être causée par 1 ou plusieurs bactéries
    \item \textbf{Exemple} : "[Bactériémie] à [E. coli] et [K. pneumoniae]"
\end{itemize}

\subsubsection{Relation ORIGINE}
\begin{itemize}
    \item \textbf{Type} : Bactériémie → Site\_primaire
    \item \textbf{Cardinalité} : 1 bactériémie a 0 ou 1 site primaire identifié
    \item \textbf{Exemple} : "[Bactériémie] sur [cathéter central]"
\end{itemize}

\subsubsection{Relation COMPLICATION}
\begin{itemize}
    \item \textbf{Type} : Bactériémie → Site\_secondaire
    \item \textbf{Cardinalité} : 1 bactériémie peut avoir 0 à N sites secondaires
    \item \textbf{Exemple} : "[Bactériémie] compliquée d'[endocardite]"
\end{itemize}

\subsubsection{Relation PROFIL\_RÉSISTANCE}
\begin{itemize}
    \item \textbf{Type} : Bactérie → Résistance
    \item \textbf{Cardinalité} : 1 bactérie a 0 ou 1 profil de résistance documenté
    \item \textbf{Exemple} : "[S. aureus] [résistant à la méthicilline]"
\end{itemize}

\subsection{Règles d'annotation spécifiques}

\subsubsection{Périmètre d'annotation}
\begin{enumerate}
    \item \textbf{Règle fondamentale} : Annoter UNIQUEMENT les éléments liés à une bactériémie confirmée
    \item Les infections sans bactériémie associée ne sont PAS annotées
    \item L'information doit être recherchée prioritairement dans la même phrase que la mention de bactériémie
    \item Extension possible au paragraphe si nécessaire pour la cohérence
\end{enumerate}

\subsubsection{Gestion de la négation}
Ne PAS annoter les éléments précédés de termes de négation :
\begin{itemize}
    \item "absence de", "pas de", "sans"
    \item "a écarté", "écarte", "exclu"
    \item "ne révèle pas", "ne montre pas"
    \item "peu probable", "douteux"
\end{itemize}

\subsubsection{Cas particuliers}
\begin{itemize}
    \item \textbf{Bactériémies multiples} : Annoter séparément chaque épisode
    \item \textbf{Co-infections} : Annoter toutes les bactéries identifiées
    \item \textbf{Sites multiples} : Distinguer site primaire unique et sites secondaires multiples
    \item \textbf{Évolution temporelle} : Annoter l'état au moment de la rédaction du CR
\end{itemize}

\newpage